\section{The quadratic relations of a general Gaussian eightfold}
\label{sec:theta-eightfolds}
This section tests the most direct extension of van Geemen's fourfold construction and identifies exactly where that extension stops.

For an even principal Gaussian abelian variety of dimension $2n$,
with the theta structure used in \cite{vG96}, van Geemen's
Corollary~2.4 gives
\[
 \dim H^0(A,2\Theta)_-=2^{n-1}(2^n-1).
\]
Here the subscript denotes the minus eigenspace of the induced
involution on second-order theta sections.
His Lemma~3.3 identifies the base locus of the minus projection on an indecomposable member as $2^{2n}$ two-torsion points \cite[pp.~621--622]{vG96}. For $n=2$, Proposition~3.6 produces a smooth quadric fourfold in the resulting $\mathbf P^5$; Theorem~3.7 then proves that the pullback of a ruling plane detects the Weil classes. For $n=4$, the same general statements give $120$ sections, a map to $\mathbf P^{119}$, and $256$ base points. They do not give a quadric eightfold in $\mathbf P^9$.

We examine the precise space of quadratic relations obtained by applying the same symmetry to the theta identities. Put $G=\mathbf F_2^4$ and let $V$ have basis $(e_u)_{u\in G}$. The second-order theta coordinates form a matrix $(X_{uv})_{u,v\in G}$. The involution transposes this matrix, so its minus eigenspace identifies with $\bigwedge^2V$. The vanishing theta constants in \cite[Lemma~2.3 and \S3.5]{vG96} give the $120$ quadratic forms
\begin{equation}\label{eq:theta-forced}
 Q_{a,b}(X)=\sum_{u,v\in G}(-1)^{(u+v)\cdot b}
 X_{uv}X_{u+a,v+a},\qquad a\cdot b=1.
\end{equation}
A linear combination descends through the minus projection when it vanishes on the plus space of symmetric matrices. There are no mixed terms, since each form is invariant under transposition.

\begin{theorem}\label{thm:theta-rank}
The combinations of \eqref{eq:theta-forced} that descend to the minus eigenspace form a vector space of dimension $35$. Every nonzero form in this space has rank at least $24$, and rank $24$ occurs. In particular this space contains no rank-$10$ quadric.
\end{theorem}

We first prove this linear-algebra assertion. We then show that these
are all the quadratic relations of the minus projection on a nonempty
Zariski-open subset of the even principal Gaussian family.

\begin{proof}
For $a\cdot b=1$ define the skew matrix
\[
 (P_{a,b})_{u,w}=(-1)^{u\cdot b}\,\mathbf1_{w=u+a}.
\]
For fixed nonzero $a$, the eight permitted characters form a basis of the skew matrices supported on displacement $a$. Different displacements have disjoint support. Consequently these $120$ matrices form a basis of $\bigwedge^2V^*$.

Associate to $Q=\sum c_{a,b}Q_{a,b}$ the symmetric tensor
$T=\sum c_{a,b}P_{a,b}\otimes P_{a,b}$ on $\bigwedge^2V$. Substitution gives
\[
 Q(vv^{\mathsf t}+ww^{\mathsf t})=2T(v\wedge w,v\wedge w).
\]
Since every complex symmetric matrix is a sum of rank-one symmetric matrices, $Q$ vanishes on the symmetric matrices if and only if $T(v\wedge w,v\wedge w)=0$ for every $v,w$. Polarizing first in $w$ and then in $v$ gives
\[
 T(x\wedge y,x\wedge z)=0,\qquad
 T(x\wedge y,w\wedge z)=-T(w\wedge y,x\wedge z).
\]
Together with the two wedge antisymmetries, these identities say that $T$ is induced by an alternating four-form $\xi\in\bigwedge^4V^*$:
\[
 T_\xi(x\wedge y,z\wedge w)=\xi(x,y,z,w).
\]
Exchanging the middle two indices shows that the restriction of $Q$ to the skew matrices is a nonzero scalar multiple of this form. Thus their ranks agree.

Let $\tau_s e_u=e_{u+s}$ and $\epsilon_t e_u=(-1)^{t\cdot u}e_u$. On the skew matrices $P_{a,b}$ these operations have characters $(-1)^{b\cdot s}$ and $(-1)^{a\cdot t}$. The pairs of characters are distinct. Hence their invariant symmetric tensors are exactly the diagonal tensors $\sum c_{a,b}P_{a,b}^2$. It follows that the descending forms identify with the four-forms invariant under every $\tau_s$ and $\epsilon_t$.

Sign invariance forces $\xi(e_u,e_v,e_w,e_z)=0$ unless $u+v+w+z=0$. For distinct indices this means that they form an affine two-plane. Translation identifies the coefficients on the four cosets of a fixed two-dimensional subspace $P\subset G$. Translation inside $P$ is a double transposition on its four vertices and preserves its wedge orientation. There is therefore one free coefficient $c_P$ for each such subspace. Their number is
\[
 \frac{(16-1)(16-2)}{(4-1)(4-2)}=35.
\]

It remains to prove the sharp rank bound. Call $P$ active when $c_P\ne0$, let $U$ be the union of the nonzero vectors in the active planes, and put $m=|U|$. For $a\in U$, let $k_a$ count the active planes through $a$. Then
\begin{equation}\label{eq:plane-incidence}
 1\le k_a\le\left\lfloor\frac{m-1}{2}\right\rfloor,
 \qquad \sum_{a\in U}k_a=3\#\{P:c_P\ne0\}.
\end{equation}
The first bound holds because the other two nonzero vectors in different planes through $a$ are disjoint pairs.

There is an orthogonal decomposition
\[
 \bigwedge^2V=\bigoplus_{a\ne0}E_a,
 \qquad E_a=\langle e_x\wedge e_{x+a}:x\in G\rangle,
 \quad \dim E_a=8.
\]
Indeed the sign characters distinguish these summands. Choose a complement $H$ to $\langle a\rangle$ in $G$. In the basis $f_x=e_x\wedge e_{x+a}$, $x\in H$, translation invariance gives the convolution matrix
\[
 T_\xi(f_x,f_y)=b_a(y-x),\qquad
 b_a(z)=\xi(e_0,e_a,e_z,e_{z+a}).
\]
Here $H\cong\mathbf F_2^3$, $b_a(0)=0$, and $b_a$ has precisely $k_a$ nonzero values. Its eight character vectors diagonalize the matrix. If $r_a$ is its rank, then $r_a$ is the size of the Fourier support of $b_a$, and the total rank is $\sum_{a\in U}r_a$.

We record the required elementary Fourier facts. If a nonzero function $b$ on $\mathbf F_2^3$ satisfies $b(0)=0$, write $k=|\operatorname{supp}b|$ and $r=|\operatorname{supp}\widehat b|$. Then
\[
\begin{array}{c|cccccc}
\text{condition}&k=1&k=2&k=3&\text{any }k&r=2&r=3\\
\hline
\text{conclusion}&r=8&r\ge4&r\ge6&r\ge2&k=4&k=6.
\end{array}
\]
For two terms, the Fourier values have the forms $c_1+c_2$ and $c_1-c_2$, each four times. For three terms, factor out one character. The remaining two character differences are independent, so the four values are $c_1\pm c_2\pm c_3$, each twice. At most one vanishes, since two vanishing values would force a coefficient to be zero. A Fourier support of size one contradicts $b(0)=0$. A Fourier support of size two gives $b=c(\chi_1-\chi_2)$, supported on four points. For a Fourier support of size three apply the three-term argument to Fourier inversion: the value at zero vanishes, so exactly one of the four patterns vanishes, giving support size six. In particular a block with $k\le5$ has rank at least four unless $k=4$, when rank two is possible.

We now bound the sum of the ranks by considering $m$. If $m=3$, one plane is active and its three blocks have rank eight, giving $24$. The value $m=4$ is impossible. If $m=5$, exactly two planes are active and meet in one nonzero vector. No third plane can be contained in their union: a vector space cannot be covered by two proper subspaces. Four blocks have $k=1$, and the remaining block has $k=2$, giving rank at least $36$.

If $6\le m\le8$, every $k_a\le3$, so every active block has rank at least four and the total is at least $24$. If $m=9$ and $k_a=4$, the four planes through $a$ pair the other eight elements of $U$ under translation by $a$. Thus $U\cup\{0\}$, of size ten, is translation invariant under $a$. Its translation stabilizer has order a power of two dividing ten, and therefore has order at most two. At most one $a$ can have $k_a=4$. The rank is at least $8\cdot4+2=34$.

For $m=10$, the baseline rank is $20$. A total below $24$ would exclude $k=1$ and $k=3$ blocks, whose excesses above two are at least six and four. It would permit at most one $k=2$ block. All other $k$'s would be four, so $\sum k_a$ would be $40$ or $38$, contrary to divisibility by three in \eqref{eq:plane-incidence}. For $m=11$, $k_a\le5$, so rank three is impossible. A total below $24$ would make every block rank two and every $k_a=4$, giving the impossible sum $44$. Finally, if $m\ge12$, the rank is at least $2m\ge24$.

To attain $24$, take the invariant orbit sum for a single two-plane $P$. Its four terms are volume forms on the four disjoint coordinate four-spaces indexed by its affine cosets. Each volume form induces a nondegenerate pairing of rank six on the exterior square of its four-space. Mixed wedges are in the radical. The resulting rank is $4\cdot6=24$.
\end{proof}

\begin{lemma}\label{lem:generic-gaussian-theta-nulls}
Let $n\ge2$ and consider the even principal Gaussian period domain
\[
 \mathcal D_n=
 \left\{\tau=
 \begin{pmatrix}X&Y\\-Y&X\end{pmatrix}\in\mathfrak H_{2n}:
 X^{\mathsf t}=X,\ Y^{\mathsf t}=-Y\right\}.
\]
On a nonempty open subset, the only vanishing even theta constants
are those whose two genus-$n$ characteristics are equal and odd.
Their number is $2^{n-1}(2^n-1)$.
\end{lemma}

\begin{proof}
Write a characteristic as a pair $(m_1,m_2)$, where
$m_j=(a_j,b_j)\in\mathbf F_2^n\times\mathbf F_2^n$ and its parity is
$a_j\cdot b_j$. An even characteristic of genus $2n$ has the same
parity in the two factors. Changing summation variables
$(u,v)\mapsto(-v,u)$ in the theta series shows
\[
 \vartheta_{m_1,m_2}(\tau,0)
   =(-1)^{a_2\cdot b_2}\vartheta_{m_2,m_1}(\tau,0).
\]
Thus the equal odd pairs vanish identically, as in
\cite[Lemma~2.3]{vG96}.

At $Y=0$ the constant factors as
$\vartheta_{m_1}(X,0)\vartheta_{m_2}(X,0)$.
Every even theta constant is a nonzero holomorphic function on
$\mathfrak H_n$: symplectic change of marking reduces it to the
zero characteristic, nonzero at a product of elliptic periods.
The product of two nonzero holomorphic functions on this connected
domain is nonzero. Therefore a pair of even characteristics does
not vanish identically on $\mathcal D_n$.

Suppose $m_1,m_2$ are distinct odd characteristics, and put
$g_m(X)=\nabla_z\vartheta_m(X,0)$.
Differentiating the convergent theta series in a skew matrix
direction $B$ gives, at $Y=0$,
\[
 d_Y\vartheta_{m_1,m_2}(X,0)[B]
   =\frac{1}{2\pi i}\,g_{m_1}(X)^{\mathsf t}B\,g_{m_2}(X).
\]
It suffices to find one $X$ where these two gradients are not
proportional.

The integral symplectic group is transitive on ordered pairs of
distinct odd characteristics. Here is the finite-dimensional reason.
Characteristics are quadratic refinements of the symplectic form
on $\mathbf F_2^{2n}$, and parity is their Arf invariant.
Two odd refinements differ by pairing with a unique nonzero vector
$v$ isotropic for the first refinement. The orthogonal group of
that odd refinement is transitive on its nonzero isotropic vectors:
complete $v$ to a hyperbolic pair, and identify the orthogonal
complements using the standard normal form of a nonsingular
quadratic form over $\mathbf F_2$. The finite symplectic changes
lift to integral ones, since symplectic transvections generate
the group and lift integrally.

We may consequently reduce to two odd characteristics with their
only odd elliptic component in the first and second coordinates,
respectively. At a product of elliptic periods their gradients
are nonzero multiples of the first and second coordinate vectors.
Indeed the odd elliptic theta function has a simple zero at the
origin and each even elliptic theta constant is nonzero.
The theta transformation formula takes an odd gradient to a
nonzero scalar times one common invertible linear transformation
of that gradient; its other differentiated term is zero because
the odd theta constant is zero. Independence of the two gradients
therefore persists after undoing the change of marking.
The displayed derivative is nonzero for a suitable skew $B$.

Every even theta constant other than an equal odd pair is thus
nonzero somewhere on the connected domain $\mathcal D_n$.
There are only finitely many such constants, so the complement
of their zero loci is nonempty and open. The algebraic form of
the conclusion also follows below from the rank of the
multiplication map.
\end{proof}

\begin{corollary}\label{cor:all-eightfold-quadrics}
On a nonempty Zariski-open subset of each marked even principal
Gaussian eightfold component, the entire quadratic ideal of the
minus-projected image in $\mathbf P^{119}$ is the
$35$-dimensional space in Theorem~\ref{thm:theta-rank}.
Every nonzero quadric in that ideal has rank at least $24$.
In particular no ten sections of this minus eigenspace can map
the general fibre dominantly to a smooth quadric eightfold.
\end{corollary}

\begin{proof}
Let $g=8$ and let $X_\sigma$, $\sigma\in\mathbf F_2^g$, be the
second-order theta coordinates. The quadratic forms
\[
 Q_{a,b}=\sum_\sigma(-1)^{\sigma\cdot b}
                 X_\sigma X_{\sigma+a},
 \qquad a\cdot b=0,
\]
form a basis of all quadratic forms: for a fixed displacement $a$,
the permitted characters form the Fourier basis on the unordered
pairs $\{\sigma,\sigma+a\}$; when $a=0$ they form the basis on
the squares. The addition formula
\cite[\S3.5, pp.~622--623]{vG96} says
\[
 Q_{a,b}\bigl(\vartheta[\sigma;0](2\tau,2z)\bigr)
 =\vartheta[a;b](\tau,0)\,\vartheta[a;b](\tau,2z).
\]
The sections $\vartheta[a;b](\tau,2z)$, as $(a,b)$ ranges over the
even characteristics, form a basis of the even sections of $4\Theta$.
It follows that the full Kummer quadratic ideal is precisely the
span of the quadrics for vanishing even theta constants.

Lemma~\ref{lem:generic-gaussian-theta-nulls} makes this kernel,
on a nonempty open set, the span of the $120$ forms
\eqref{eq:theta-forced}. A quadric on the minus projection
pulls back to a polynomial involving only minus coordinates.
Its vanishing is equivalent to membership in the full kernel.
The intersection of that kernel with these polynomials is
exactly the $35$-dimensional space already computed.
Conversely every such polynomial vanishes on the projected image.

After a finite level cover the spaces of sections and their
multiplication are algebraic vector bundles and an algebraic
bundle map. The minimal kernel dimension is attained on a
Zariski-open set. Thus the nonempty analytic open just found
gives the asserted nonempty Zariski-open set.
Ten independent sections defining a map to a smooth $Q^8$
would pull its equation back to a nonzero rank-$10$ quadratic
form on the minus space, contradicting the rank bound.
\end{proof}

The smallest nonzero relation on this open set has a smooth
quotient quadric of dimension $22$. Its degree-eight cohomology
is below the middle and is generated by a hyperplane power.
Pulling back that class gives a multiple of $\theta^4$, whose
Weil projection is zero. Thus the direct ruling-class argument
cannot use any quadric of the general minus projection.
Special fibres may have additional relations. Cycles involving
singular loci, sections outside the minus space, and different
target varieties require other arguments.

\begin{proposition}[What a quadric eightfold would give]\label{prop:quadric-gate}
Let $(A,L)$ be a Weil eightfold and $\Delta=\chi(L)$. Suppose ten sections of $2L$ define a dominant rational map to a smooth quadric $Q^8\subset\mathbf P^9$, with base locus of dimension at most three. Suppose that, for a general member $P\cong\mathbf P^4$ of one ruling, the closure of its inverse image has nonzero Weil projection. Then
\[
 m_W(A,L)\le2,\qquad \delta_W(A,L)\le8\cdot8!\Delta=322560\Delta.
\]
\end{proposition}

\begin{proof}
A general $\mathbf P^5$ containing a ruling $P$ cuts $Q^8$ in $P\cup P'$, where $P'$ is a member of the other ruling. The $\mathbf P^5$ is cut out by four hyperplanes in $\mathbf P^9$. By transitivity of the orthogonal group and the dimension theorem for a general translate, the inverse images of a general such pair of ruling planes have dimension four away from the base locus. The base locus has dimension at most three. Thus the four corresponding members of $|2L|$ have proper intersection. Their intersection contains the closure of the inverse image of $P$, with positive multiplicities. Its nonzero Weil projection implies that one irreducible component detects. Proposition~\ref{prop:complete-intersections} gives
$\delta_W\le(8!\Delta/2)2^4=8\cdot8!\Delta$.
\end{proof}

The geometric hypotheses of this proposition remain unproved in
the Gaussian family. Corollary~\ref{cor:all-eightfold-quadrics}
excludes obtaining its rank-$10$ quadric from any ten sections
of the general minus eigenspace.
